%% kapa: exam=1 round=33 year=2022 period=1 form=A subject=civil_law subject_ko=민법 seq=26 number=26
%% source: Q-Net 감정평가사 기출문제 (cst00309) artlSeq=5209997 file=raw/1008/5209997/2022년도 제33회 감정평가사 1차시험 1교시(민법,경제학,부동산학).pdf
%% conversion: kordoc 4.24.51 -> tools/gen1.py (auto); review=proofread-sonnet; flags=none
\begin{question}{26}
주위토지통행권에 관한 설명으로 옳지 않은 것은? (다툼이 있으면 판례에 따름)
\begin{choices}
  \item 토지의 분할로 주위토지통행권이 인정되는 경우, 통행권자는 분할당사자인 통행지 소유자의 손해를 보상하여야 한다.
  \item 통행지 소유자는 통행지를 배타적으로 점유하고 있는 주위토지통행권자에 대해 통행지의 인도를 청구할 수 있다.
  \item 주위토지통행권은 법정의 요건을 충족하면 당연히 성립하고 요건이 없어지면 당연히 소멸한다.
  \item 주위토지통행권에 기한 통행에 방해가 되는 축조물을 설치한 통행지 소유자는 그 철거의무를 부담한다.
  \item 주위토지통행권의 범위는 현재의 토지의 용법에 따른 이용의 범위에서 인정된다.
\end{choices}
\end{question}
\answer{1}
